
\chapter{Literature Review}
\label{ch:litreview}

\section{Introduction}
\label{sec:lit_intro}

Plastic waste sorting faces two coupled challenges: a global recycling shortfall and a long-standing robotic automation bottleneck. Globally, less than 10\% of plastic waste is appropriately recycled, with a large proportion lost to inadequate sorting infrastructure~\citep{ramos2024}. Industrial Material Recovery Facilities (MRFs) still depend heavily on near-infrared (NIR) optical sensors and rigid rule-based controllers that struggle with material variability, deformation, and contamination~\citep{lubongo2024, kroell2022}. Vision-Language-Action (VLA) models offer a different premise: by integrating visual perception, language instruction, and action generation in a single architecture pre-trained on internet-scale data, they promise generalisation across the diverse morphologies of real waste without hard-coded rules per plastic type~\citep{sapkota2025, kawaharazuka2025}.

Translating that promise to sorting lines, however, collides with several unsolved problems, two of which this dissertation isolates: transparent plastics defeat the depth sensors on which manipulation pipelines rely, and no corpus of robot demonstrations exists on which to fine-tune a VLA for this domain. This chapter reviews the classical sorting literature (Section~\ref{sec:lit_classical}), scene understanding for robotic sorting with emphasis on transparent-object perception (Section~\ref{sec:lit_scene}), the VLA paradigm and its documented spatial limitations (Sections~\ref{sec:lit_vla}--\ref{sec:lit_spatial}), and synthetic-data strategies for training such systems (Section~\ref{sec:lit_synthetic}). Section~\ref{sec:lit_gaps} consolidates the technical gaps, and Section~\ref{sec:lit_aim} states the research aim and questions that follow from them.

\section{Classical Plastic Sorting: Sensing and Classification}
\label{sec:lit_classical}

\citet{lubongo2024} survey current MRF sorting technology, documenting three broad approaches: mechanical separation by density, electromagnetic extraction of ferrous materials, and optical sorting using spectral sensors. Among these, NIR spectroscopy has become the industrial standard: it permits contact-free identification of polymer types from their molecular absorption spectra, and commercial NIR systems achieve throughputs of 60--80 picks per minute with classification accuracies above 99\% under controlled conditions~\citep{lubongo2024}.

Classical systems nonetheless face structural limitations. Standard NIR equipment cannot identify black plastics: the carbon-black pigment absorbs nearly all incident light across the 0.9--1.7\,$\mu$m range, rendering these items invisible to the sorter~\citep{kroell2022}. Mid-wave infrared hyperspectral imaging (3--5\,$\mu$m), where plastics exhibit distinct absorption features regardless of colour, has emerged as the most effective response~\citep{kroell2022}, at considerably higher hardware and processing cost. Contamination is a second failure mode: food residue and labels degrade the accuracy of both NIR and visual sensors~\citep{kroell2022, ramos2024}, and rule-based systems lack the contextual reasoning to recognise a contaminated PET bottle as PET.

On the machine-vision side, RGB-based classification has progressed steadily. \citet{bobulski2021} benchmark CNNs for polymer classification from low-cost RGB cameras, reporting 74\% cross-validated accuracy for a lightweight 15-layer network and 86\% for MobileNet-v2---sufficient for broad material categories but not fine-grained polymer discrimination. One-stage detectors such as YOLO variants have become the dominant choice for real-time conveyor detection~\citep{wu2023}. \citet{kroell2022} conclude from their systematic review that while deep learning has improved classification flexibility, distinguishing chemically distinct but visually identical polymers---clear PET from clear PVC in particular---remains unsolved by RGB methods alone, making spectral data a necessity wherever chemical accuracy is required~\citep{lubongo2024, yang2022}.

A further, less discussed limitation is geometric: spectroscopic and RGB classification answer \emph{what} is on the belt but not \emph{where to grasp it}. Converting a detection into a successful pick on cluttered, deformed waste is repeatedly identified as harder than the classification itself~\citep{lubongo2024}. \citet{koskinopoulou2021} present an integrated vision-based robotic sorting system on an industrial line---combining deep-learning categorisation, grasp selection, and robotic separation, with reported sorting accuracy above 90\%---which exemplifies the state of practice and, equally, its constraints: the perception stack presumes usable depth sensing, and extending the system to new categories or new sorting instructions requires retraining and re-engineering rather than issuing a new command.

\section{Semantic Scene Understanding for Robotic Sorting}
\label{sec:lit_scene}

Transforming raw sensor data into actionable representations for sorting rests on three components: open-vocabulary detection, depth reconstruction, and affordance prediction.

Traditional detectors require a fixed category vocabulary defined at training time, making them brittle against novel packaging. VoxPoser~\citep{huang2023} demonstrates a more flexible alternative by chaining an open-vocabulary detector (OWL-ViT) with the Segment Anything Model, identifying and masking arbitrary objects described in natural language---``the crushed translucent bottle''---without labelled examples of that object. This capability motivates language-conditioned sorting, but VoxPoser's downstream machinery composes LLM-generated 3D value maps over voxel grids and therefore \emph{presumes reliable 3D perception}---precisely what transparent objects deny.

Depth reconstruction for transparent objects is the second component. Active infrared sensors such as the Intel RealSense series project structured light that refracts through or transmits entirely across transparent PET and PVC, producing depth maps with severe holes and distortions exactly over the objects most common in plastic waste~\citep{sajjan2020}. ClearGrasp~\citep{sajjan2020} addresses this directly: from the RGB image it predicts surface normals, transparent-object masks, and occlusion boundaries via three parallel DeepLabV3-based networks, then solves a global optimisation that reconstructs depth over the corrupted regions, anchored to the reliable measurements of the surrounding opaque scene. Trained purely on a 45{,}000-image physically based synthetic dataset (which also supplies the assets used in the present work), ClearGrasp raised parallel-jaw grasp success on real transparent objects from a 12\% raw-depth baseline to 74\%, establishing that the RGB-to-geometry route bridges the sim-to-real gap for this material class. Two properties of ClearGrasp are inherited here---the decomposition into RGB-based geometric prediction plus optimisation-based completion, and synthetic supervision---while one is deliberately departed from: ClearGrasp terminates in a task-specific grasp network with no semantic or language interface.

Affordance prediction closes the loop from perception to action. VoxPoser's LLM-composed value maps assign reward to actionable regions such as a bottle neck~\citep{huang2023}. RoboPoint~\citep{yuan2024} instead trains a vision-language model to predict 2D spatial coordinates for grasp targets directly, trading 3D resolution for inference efficiency---an approach structurally close to the coordinate-token prediction evaluated in this dissertation. SpatialVLA~\citep{qu2025} extends the idea by injecting egocentric 3D representations into the VLA action-token space to better resolve spatial relations in clutter.

\section{Vision-Language-Action Models}
\label{sec:lit_vla}

The VLA paradigm restructures robotic system design: rather than composing separate modules for perception, planning, and control, a single model learns an end-to-end mapping from visual observations and language instructions to actions~\citep{sapkota2025, kawaharazuka2025}. Surveys catalogue over eighty VLA models published since 2022, categorised by action-generation strategy---autoregressive token prediction, diffusion, and flow matching~\citep{ma2024, li2025survey, zhong2025}---and by functional role: sensorimotor models, world models, and affordance-based models, the latter particularly suited to unstructured environments because the generated action is decoupled from a specific robot's kinematics~\citep{kawaharazuka2025}.

RT-2~\citep{zitkovich2023} established the token approach at scale: continuous end-effector displacements and gripper commands are discretised into 256 bins and generated as text, and co-fine-tuning on web-scale vision-language data alongside robot trajectories prevents catastrophic forgetting of semantic knowledge while acquiring physical skill. RT-2's headline capability---manipulating objects never seen in the robot data by drawing on web knowledge---is precisely the property that makes VLAs attractive for waste streams of unbounded variety. RT-2 is, however, closed and computationally heavy (its 55B variant runs at 1--3\,Hz on multi-TPU cloud hardware; the 5B variant at $\sim$5\,Hz~\citep{zitkovich2023}).

OpenVLA~\citep{kim2024} made the paradigm accessible: a 7B-parameter open-source model pairing a fused SigLIP--DINOv2 vision encoder with a Llama-2 backbone, trained on 970k real-robot episodes, adopting the same 256-bin action tokenisation. Critically for resource-constrained settings, \citet{kim2024} demonstrate effective adaptation with Low-Rank Adaptation~\citep{hu2021lora} from as few as tens of demonstrations, and quantised operation~\citep{dettmers2023qlora} brings fine-tuning within reach of a single high-memory GPU---the configuration adopted in Chapter~3. The choice of tokenisation itself is an active research area~\citep{zhong2025}; uniform binning is the simplest scheme, and its 256-level granularity bounds achievable spatial precision. Complementary directions add intermediate reasoning (visual chain-of-thought in CoT-VLA~\citep{zhao2025}), explicit 3D world models~\citep{zhen2024}, or additional physical sensing modalities (OmniVLA~\citep{omnivla2025}).

\section{Spatial Grounding: a Documented Weakness}
\label{sec:lit_spatial}

A consistent finding across recent evaluations is that VLAs, and the vision-language models beneath them, are semantically strong but metrically weak. SpatialBench~\citep{cai2024} documents systematic deficits in localisation and spatial-relation reasoning even where recognition is nearly perfect. The affordance and representation work of Section~\ref{sec:lit_scene}---RoboPoint's dedicated pointing supervision~\citep{yuan2024}, SpatialVLA's 3D encodings~\citep{qu2025}, 3D-VLA's generative world model~\citep{zhen2024}---can be read collectively as responses to this weakness: spatial grounding does not emerge from scale alone but must be engineered through data, representation, or architecture. For sorting, the implication is sharp: recognising a bottle is insufficient; the policy must emit coordinates accurate to a gripper's tolerance. This dissertation contributes evidence from an unusual angle---measuring how far purely synthetic, composited 2D supervision pushes spatial grounding in a stock OpenVLA, and characterising what happens when the data is insufficient.

\section{Synthetic Data and Its Hazards}
\label{sec:lit_synthetic}

Collecting real robot demonstrations is expensive; for transparent waste no public corpus exists. Two synthetic strategies dominate. Physically based rendering, as in ClearGrasp~\citep{sajjan2020}, yields high-fidelity images with exact ground truth at the cost of 3D assets and rendering infrastructure. Compositing---pasting segmented foregrounds onto varied backgrounds, as in the ReSort-IT augmentation methodology of \citet{koskinopoulou2021} adopted and extended in this work---is far cheaper, at the cost of physical implausibility: composites lack correct shadows and contact regions and, for transparent objects, refraction of the true background through the object. Domain randomisation of pose, scale, photometry, and backgrounds is the standard mitigation, forcing reliance on object-intrinsic cues.

The literature is comparatively quiet on a hazard specific to composited \emph{action} supervision: if target labels correlate with background identity, or compositing artefacts such as sharp alpha seams are visible, a high-capacity model can minimise training loss through shortcuts carrying no spatial understanding---and, because most dimensions of an action string are near-constant, aggregate token metrics can conceal the failure entirely. Chapter~4 documents a clean instance of this failure and its diagnosis, contributing an empirical data point on what composited data must control for when the supervised quantity is a spatial action rather than a class label.

\section{Technical Gaps for VLA-Based Plastic Sorting}
\label{sec:lit_gaps}

Synthesising the review, four structural mismatches separate current VLA systems from deployment in plastic sorting.

\textbf{Gap 1: Inference latency.} Autoregressive VLAs generate actions token by token and cannot reach the control rates of reactive manipulation: RT-2 runs at 1--5\,Hz depending on scale~\citep{zitkovich2023}, OpenVLA at roughly 5\,Hz on a high-end GPU~\citep{kim2024}, and \citet{kawaharazuka2025} identify this as the central deployment barrier. Sorting lines compound the problem with continuously moving targets.

\textbf{Gap 2: The embodiment and domain gap.} VLA training data consists overwhelmingly of clean, well-lit household manipulation on specific arms; industrial waste---deformed, contaminated, cluttered, and moving---is out of distribution both visually and kinematically~\citep{kawaharazuka2025}. Classical pipelines such as \citet{koskinopoulou2021} remain the benchmark on real waste precisely because they were engineered in-domain.

\textbf{Gap 3: Data scarcity for transparent waste.} No demonstration corpus exists for transparent-waste manipulation, and transparent objects additionally invalidate the depth channel that RGB-D policies and 3D affordance methods presume (Section~\ref{sec:lit_scene}). Synthetic generation is the only accessible route, and Section~\ref{sec:lit_synthetic} identified its under-documented shortcut hazards for action supervision.

\textbf{Gap 4: Absence of spectral integration.} Current VLAs operate on RGB or RGB-D inputs and cannot chemically discriminate visually identical polymers such as clear PET and clear PVC~\citep{lubongo2024}; OmniVLA's multi-sensor representations~\citep{omnivla2025} have not been applied to NIR polymer classification. This dissertation operates deliberately \emph{within} that constraint: its transparent target categories are drawn from the ClearGrasp object set and are distinguished morphologically, chemical identification being excluded from scope and identified as future work.

Table~\ref{tab:lit_summary} summarises the principal approaches against these gaps.

\begin{table}[htbp]
\centering
\caption{Summary of key approaches in automated plastic sorting and robotic manipulation.}
\label{tab:lit_summary}
\renewcommand{\arraystretch}{1.3}
\begin{tabular}{@{}p{2.8cm}p{2.3cm}p{3.2cm}p{6.2cm}@{}}
\toprule
\textbf{Study} & \textbf{Modality} & \textbf{Architecture} & \textbf{Key contribution and limitation} \\
\midrule
\citet{lubongo2024} & NIR spectroscopy & Rule-based & Industrial standard, $\sim$99\% accuracy under controlled conditions; blind to black plastics; no grasp geometry. \\
\citet{kroell2022} & RGB vision & CNN / YOLO & Flexible visual classification; cannot separate chemically distinct, visually identical polymers (PET vs.\ PVC). \\
\citet{koskinopoulou2021} & RGB(-D) & Deep classical pipeline & Integrated robotic sorting on an industrial line; presumes valid depth; retraining required per category. \\
\citet{sajjan2020} & RGB-D & DeepLabV3 + optimisation & Transparent depth reconstruction (grasp success 12\%$\to$74\%); no semantic or language interface. \\
\citet{zitkovich2023} & RGB + text & VLA (co-fine-tuned) & Web-knowledge transfer to control; closed; 1--5\,Hz inference. \\
\citet{kim2024} & RGB + text & Prismatic VLM + Llama-2 & Open-source VLA; LoRA-adaptable on one GPU; inherits latency and spatial-precision limits. \\
\citet{omnivla2025} & Multi-sensor & Sensor-fused VLA & Non-RGB modalities in VLA representations; unvalidated for NIR sorting. \\
\bottomrule
\end{tabular}
\end{table}

\section{Research Gap, Aim, and Questions}
\label{sec:lit_aim}

No prior work, to the author's knowledge, combines the elements this review has assembled: a VLA fine-tuned for transparent-waste grasp targeting using exclusively composited synthetic supervision, evaluated jointly with a learned depth-repair front-end, with the error propagation between the two quantified and the shortcut hazards of composited action data documented. Within the resource envelope of an MSc project---one GPU, no demonstration corpus, no guaranteed robot access---the deployable-system question of the original proposal is deliberately narrowed to the prior question on which it depends: \emph{under what data conditions does synthetic-only VLA adaptation produce genuine spatial grounding on transparent objects, and how do the perception and policy errors compose?}

The aim of this work is therefore to design, implement, and quantitatively evaluate a cascaded pipeline for language-conditioned grasp-target prediction on transparent plastic waste, trained entirely on synthetic data, answering:

\begin{itemize}
    \item[\textbf{RQ1}] Can procedurally composited synthetic data provide sufficient supervision for a VLA to learn image-conditioned spatial grounding of transparent objects, and what properties of the data determine success or failure?
    \item[\textbf{RQ2}] How accurately does an RGB-only network recover the surface geometry of transparent objects, and how does that accuracy degrade on unseen object categories?
    \item[\textbf{RQ3}] Does surface-normal-guided depth reconstruction improve on geometry-unaware inpainting, and how does error in the predicted normals propagate into reconstructed depth?
    \item[\textbf{RQ4}] Is the integrated pipeline computationally feasible for conveyor deployment, and if not, which component bounds its throughput?
\end{itemize}

\textbf{[PLACEHOLDER --- add RQ5 here once the semantic-reasoning experiment is under way: Does the fine-tuned policy condition its spatial predictions on the language instruction, selecting between different object categories present in the same scene?]}

The following chapter details the methodology through which these questions are answered.